\chapter{An interior small-$L$ C1 source: a full three-circuit closure}
\label{component:pe3ef590c01d6}
\noindent\textbf{Component }\path{pe3ef590c01d6}\par
{\small\noindent Source: \path{R13/advances/c1_small_interior/interior_closure.tex}\par}
{\small\noindent SHA256: \path{e3ef590c01d68ef36079d3024a86c5d75167e0b3ae52d689edad8ce6b1da744b}\par}
{\small\noindent Inventory use: IC7--IC8, IC11--IC19 and finite arrays | IC11-IC21 only | C14\_\allowbreak{}ENTIRE\_\allowbreak{}ORIGINAL\_\allowbreak{}PARENT\_\allowbreak{}GRAPH; precise subordinate uses retained below\par}
\medskip
\section[An interior small-L C1 source: a full three-circuit closure]{An interior small-$L$ C1 source: a full three-circuit closure}
\label{pe3ef590c01d6:sec:r9-c1-small-interior}

\begin{theorem}
Use the original physical native C1 source and all earlier intersections.
Suppose $x=k^2$ is infinitesimal, and
\[
 \eta,\ \zeta,\ r,\ 1-r,\ j_0=\eta+r-1
       \quad\hbox{are positive bounded units}.                 \tag{IC1}
\]
Then the full original parent is tensor SOS for every independent
projective test and every state. No source/testing rate comparison is
used. This is a subdomain of C14.3; it does not assert the remaining
small-$L$ source boundaries or C14.4.
\end{theorem}

Normalize $d=1$, set $\rho=\eta+r$, $y=x\eta\zeta$, and retain
all original packets. Here $L=x\zeta$ is infinitesimal. The source
formulas (NC1)--(NC4) and (SB1)--(SB4) give
\[
 \begin{gathered}
 m/x\sim\rho/(2\zeta),\qquad
 \alpha/x^2\sim\rho\zeta(1+\zeta)/2,\\
 C=C_0+O(x),\quad A=A_0+O(x),\quad D=D_0+O(x),\\
 C_0={\rho\over2(1-r)},\quad A_0={\eta\over2r},
 \quad D_0={1\over2j_0}.                                \end{gathered} \tag{IC2}
\]
If the total scale is a unit, the earlier positive-total-scale theorem
already applies. On the residual total-scale-infinitesimal intersection,
all original coefficients except $\beta$ are bounded, so $\beta$ is
unbounded. Its exact numerator implies
\[
 \theta=x\pi,\quad\pi\in\operatorname{Inf}_{>0},\qquad
 \beta\theta/x=O(1),\quad
 \beta\theta^2/x^2=O(\pi).                               \tag{IC3}
\]
Thus $\beta=0$ and $D=0$ are absent on this particular residual source
subdomain, by source implication. No permitted boundary is removed
from a larger claimed theorem. All source denominators used below are
products of positive bounded units from (IC1), $M,J,Q,P_d$, and their
source specializations. The apparent $x$ and $x^2$ denominators cancel
by finite polynomial numerator division, as specified below.

The existing theorem in Section~\ref{p6d2f659c0fde:sec:r8-c1-small-L-bounded-eta},
equations (SB6)--(SB12), covers the original source $x,L$ infinitesimal,
$\zeta$ noninfinitesimal, and $\eta$ bounded on the complement of
$u,v$ both infinitesimal, with every homogeneous endpoint. Its source
domain contains (IC1). Literal testing axes are direct: the half core
is $D_v(X^2+Y^2)/2$ at $u=0$ and
$\{(uW-Y)^2+(X+uZ)^2\}/2$ at $v=0$, and every other packet is
nonnegative. The exact transformations
$(u,Y,Z)\mapsto(-u,-Y,-Z)$ and
$(v,X,Z)\mapsto(-v,-X,-Z)$ preserve the original parent, so work on
the positive testing sign chart and restore all signs afterwards.
For complex states apply the real identity to the real and imaginary
parts. First work at $Z\ne0$, divide the state by $Z$, and set $Z=1$;
no bound on $W$ is imposed. Afterward multiply the exact identity by
the original $Z^2$. Its state-normal roots become the homogeneous
linear forms displayed below and its scalar circuit roots become
$Z$ times the roots at $Z=1$. No $Z$ inverse remains. The resulting
identity is polynomial in the original state coordinates and hence
also holds at $Z=0$.

\paragraph{A bounded-ratio frame retaining all three state normals.}
First suppose the pure-testing ratio $t=v/u$ is bounded, possibly
infinitesimal. Define
\[
 \begin{gathered}
 E=\eta uW-v,\quad
 b={r(1-r)\over\eta(\eta+2r-1)},\quad
 N=Y-ruW-bE,\\
 \xi=X+hu,\qquad N_x=\sqrt\beta\,\xi,\quad h=1-\theta.
                                                               \end{gathered} \tag{IC4}
\]
The state change is exact and uses only the pure-testing inverse $u$:
$W=(v+E)/(\eta u)$, $Y=r(v+E)/\eta+bE+N$, $X=-hu+\xi$.
Every normal vanishes at the actual source contact. Its full hard
quadratic has the strictly positive diagonal residue
\[
 (E,N,N_x):\qquad
 \operatorname{diag}\left(
 {\rho\over2\eta(\eta+2r-1)},
 {\eta(\eta+2r-1)\over2r(1-r)j_0},1\right).             \tag{IC5}
\]
This is a statement about the complete original hard block. At bounded
$t$, its remaining terms are bounded with zero residue: the potentially
new $m t^2$ coefficient is infinitesimal, and normalizing $\xi$ makes
all its other hard entries infinitesimal. No state normal has been
restricted or dropped.

\paragraph{The exact scalar, with the negative defect absorbed.}
Set $U=u^2,V=v^2$, $s=\sqrt{\eta\zeta}$, $k=\sqrt x$, and introduce
the positive separated roots
\[
 A_t=uv,\quad B_t={k v^2\over s u},\quad
 C_t=xs,\quad G=kv,\qquad F_v=V-y.
                                                               \tag{IC6}
\]
They satisfy $A_tB_tC_t=G^3$ and coincide at actual contact
$(u,t)=(k,s)$. The pure-testing pole in $B_t$ is inside the root;
none of the coefficient arrays below has a negative power of $t$.

Use $A_w,A_z,b_0,b_1$ from (SN1)--(SN2), including the beta term.
The scalar at $E=N=\xi=0$ is exactly
\[
 \mathcal S={VA_w\over\eta^2U}+A_z
                  -{2V\over\eta}(b_0+b_1U)
       =c_FF_v^2+f_AA_t^2+f_BB_t^2+f_CC_t^2-f_GG^2.    \tag{IC7}
\]
Here the complete coefficients are
\[
\begin{aligned}
 c_F={}&{C(1-r)^2+Ar^2+(\alpha+\tfrac12)r^2+2mr\over\eta^2},\\
 a_{UV}={}&{Ar^2\over\eta^2}+{h^2\over2}+2Dh+\beta\theta^2
                                     -{2b_1\over\eta},\\
 a_{UUV}={}&h^2(D+\tfrac12),\qquad
 a_{UVV}={r^2(m+A+\alpha+\tfrac12)\over\eta^2},\\
 q_V={}&{C(1-r)^2+Ar^2+(D+\tfrac12)(1-r)^2\over\eta^2}
                                     +D-{2b_0\over\eta},\\
 q_U={}&\alpha+(m+\beta+\tfrac12)\theta^2,\\
 f_A={}&a_{UV}+a_{UUV}U+a_{UVV}V,\qquad
 f_B={m\zeta\over\eta x},\\
 f_C={}&{\alpha-c_Fy^2+q_UU\over\eta\zeta x^2},\qquad
 f_G=-{q_V+2c_Fy\over x}.                              \end{aligned} \tag{IC8}
\]
Thus the source-negative quadratic base has already been combined
with the original $m$, $\alpha$, and higher-angle terms. The residues
are
\[
 c_F\sim{\rho\over2\eta^2},\qquad
 f_A,f_B,f_C\sim\lambda={\rho\over2\eta},\qquad
 f_G\sim3\lambda.                                    \tag{IC9}
\]
For example the leading scalar is exactly
\[
 {\rho\over2\eta^2}(V-y)^2+
 {\rho\over2\eta}\left\{UV+{xV^2\over\eta\zeta U}
                            +\eta\zeta x^2-3xV\right\}.
                                                               \tag{IC10}
\]
The braces are the three-term AM--GM circuit
$A_t^2+B_t^2+C_t^2-3G^2$, not a standalone quadratic PSD assertion.

\paragraph{Finite scalar Taylor arrays.}
Let $\epsilon_i=f_i-\lambda$ and $\epsilon_G=f_G-3\lambda$,
$h_c=\epsilon_A+\epsilon_B+\epsilon_C-\epsilon_G$.
Equations (IC8) show that $h_c=h_{00}+aU+cV$ with bounded
pure-source $a,c$. Exact contact value gives
\[
 h_c=a(u^2-k^2)+c(u^2t^2-k^2s^2).                     \tag{IC11}
\]
For $d_u=u-k,d_t=t-s$, its finite second Taylor identity is
\[
\begin{split}
 h_c={}&2k(a+cs^2)d_u+2k^2sc\,d_t
 +(a+ct^2)d_u^2+2kc(t+s)d_ud_t+k^2c\,d_t^2.            \end{split} \tag{IC12}
\]
The coefficient of $d_t^2$ is infinitesimal. For any affine function
$f=f_{00}+a_fU+c_fV$ among (IC8), the exact first difference is
\[
 f-f(k,s)=(u-k)(u+k)(a_f+c_ft^2)
                    +(t-s)k^2(t+s)c_f.                \tag{IC13}
\]
These are the required bounded $u$ arrays and infinitesimal $t$ arrays;
no appeal to unspecified Taylor membership is needed.
The two exact first-jet balances are
\[
 2(\epsilon_{A0}-\epsilon_{C0})+k h_{c,u0}=0,
 \quad 2(\epsilon_{B0}-\epsilon_{C0})+s h_{c,t0}=0.      \tag{IC14}
\]
They follow from the original stationary contact, and are also checked
directly using (NC4). The extra square $c_FF_v^2$ has zero value and
first jets there. Apply the accompanying three-circuit contact-lift
lemma with the exact conductors
\[
 G(u-k)=k(A_t-G),\qquad G(t-s)=s(B_t-G).                \tag{IC15}
\]
This gives a strictly positive residue Gram on the finite circuit
root list, together with the strict $F_v$ pivot.

\paragraph{Explicit finite allocations of every first normal.}
Use $f_w,f_z,k_w,k_z$ in (SN2), and define the polynomials
\[
\begin{aligned}
 P_Y(U,V)&=2(k_w/\eta+k_z),\\
 P_E(U,V)&={2\over\eta^2}(A_w/U-mV/U)
              -{2\over\eta}(b_0+b_1U)+bP_Y(U,V).
\end{aligned}                                                   \tag{IC16}
\]
The quotient $A_w/U-mV/U$ is polynomial, as its displayed numerator
shows. For either polynomial $P$ use the exact divisions
\[
 P=P_{00}+UP_U+VP_V,\quad
 P_U={P-P(0,V)\over U},\quad
 P_V={P(0,V)-P(0,0)\over V}.                           \tag{IC17}
\]
All arrays are bounded. The source formulas give, by finite numerator
division, $P_{Y,00}/x,P_{E,00}/x$ bounded. The complete first normals
of the original $F/4$ in the frame (IC4) are exactly
\[
\begin{aligned}
 \mathcal L_N&=vP_Y,\\
 \mathcal L_E&=2mv t^2/\eta^2+vP_E,\\
 \mathcal L_x&={1\over\sqrt\beta}\left[
 2u\theta(m+\beta+\tfrac12)
 -{2m\over\eta}ut^2(1+rU)+2uVZ_X\right],\\
 Z_X&=(\beta+\tfrac12)\theta-(D+\tfrac12)(1+hU)
                         +(D+\tfrac12)(1-r)/\eta .     \end{aligned} \tag{IC18}
\]
These identities are independently checked by differentiating the
entire original parent after the full state substitution (IC4).

Their actual root allocations are
\[
\begin{aligned}
 vP&={P_{00}\over k}G+u(P_U+t^2P_V)A_t,\\
 {2m\over\eta^2}vt^2&={2ms t\over\eta^2k}B_t,\\
 \mathcal L_x&={2u\theta(m+\beta+\tfrac12)
                    \over\sqrt\beta\,xs}C_t
 -{2ms(1+rU)\over\eta k\sqrt\beta}B_t
 +{2vZ_X\over\sqrt\beta}A_t.                          \end{aligned} \tag{IC19}
\]
Every coefficient in (IC19) is infinitesimal. Its $u$ first-difference
array is bounded and its $t$ array is infinitesimal: the coefficient
arrays are polynomials in $u,t$ with bounded source arrays, and each
term depending on $t$ contains $u$, $k$, or $1/\sqrt\beta$; evaluation
at $u=k$ retains such an infinitesimal factor. There are no negative
$t$ powers. The source ratios needed here are precisely
\[
 m/x,\quad P_{Y,00}/x,\quad P_{E,00}/x,\quad
 \theta(m+\beta+\tfrac12)/x,\quad \beta\theta/x,
                                                               \tag{IC20}
\]
all bounded by canceled finite numerator formulas.
The exact original contact makes each complete normal zero at
$(u,t)=(k,s)$. For a normal $\ell=\sum b_iA_i+b_GG$, put
$B=\sum b_i+b_G$, and use the explicit monic divisions
\[
 B=(u-k){B(u,t)-B(k,t)\over u-k}
       +(t-s){B(k,t)-B(k,s)\over t-s}.                 \tag{IC21}
\]
The first quotient is bounded and the second infinitesimal by the
polynomial arrays just given. Replacing $A_i$ by $A_i-G$ and applying
(IC15) allocates each complete normal with infinitesimal coefficients
on the same strict circuit roots. Joining (IC5), (IC7), and (IC19)
therefore gives a bounded positive-residue Gram for the entire original
parent in every bounded-$v/u$ testing chart, including infinitesimal
$v/u$ and arbitrary state size.

\paragraph{The reciprocal testing-ratio chart.}
It remains to consider $v/u$ unbounded, so $u/v$ is infinitesimal.
Return to homogeneous state coordinates and set
\[
 p=uW,\ q=vZ,\ E=\eta p-q,\quad
 Y=rp+bE+N,\quad B=X+uZ=\xi+\theta uZ.
\]
Use the complete raw root list
\[
 E,\ N,\ N_x=\sqrt\beta\xi,\quad
 P=\sqrt m\,vW,\quad Q=\sqrt\alpha Z,\quad
 R_t=u^2W,\quad S_t=uvW.                              \tag{IC22}
\]
No independence of these roots is required.
The exact source comparison error in the quadratic base is a
quadratic with $O(x)$ source coefficients in $(p,Y,q)$. It has zero
residue on (IC22), since
\[
 \sqrt x\,p=\sqrt{x/m}\,(u/v)P,\quad
 Y=rp+bE+N,\quad q=\eta p-E,                          \tag{IC23}
\]
and $x/m$ is a positive bounded unit. The coefficient $u/v$ is
pure testing, not a mixed source/testing comparison. The comparison
base itself is exactly $H_0N^2+\gamma E^2$, with the positive values
in (IC5).

For completeness, with $B=0$ the exact higher-angle part of the
$C,A,D$ packets and half core, after the quadratic base is removed, is
\[
\begin{split}
 &U\{AY^2-2Dq(p-Y-q)+q^2/2-pq\}\\
 &\quad+V\{C(p-Y)^2+(A+\tfrac12)Y^2\}
 +(A+\tfrac12)UVY^2+(D+\tfrac12)U^2q^2.               \end{split} \tag{IC24}
\]
Upon inserting $Y=rp+bE+N$, $q=\eta p-E$, its residue on
$(R_t,S_t)$ is strictly diagonal:
\[
 \kappa_R={\rho\eta\over2},\qquad
 \kappa_S={\rho\over2}.                              \tag{IC25}
\]
All other entries carry $u$, $v$, $U$, or $V$. For example
$UpE=uR_tE$, $VpE=vS_tE$, $U^2p^2=UR_t^2$, and
$UVp^2=US_t^2$. These identities allocate every term of (IC24)
without an inverse or a rate comparison.

The terms involving $B$ are equally explicit. Relative to $B=0$, the
$D$ packet adds $2DuvB(p-Y-D_uq)+DUVB^2$, and the half core adds
$uvB(p-Y-D_uq)+(D_v+UV)B^2/2$.
Here
\[
 B={N_x\over\sqrt\beta}+{\theta u\over\sqrt\alpha}Q,
 \quad uvp=uS_t,
\]
so their Gram coefficients are infinitesimal. The three remaining
original packets keep their entire roots:
\[
\begin{aligned}
 \sqrt m\,M&=P+\sqrt m\,uvY-\sqrt m\,B,\\
 \sqrt\alpha\sqrt{D_u}(vY-Z)
  &=\sqrt{D_u}\{\sqrt\alpha(rS_t+v(bE+N))-Q\},\\
 \sqrt\beta\sqrt{D_v}B
  &=\sqrt{D_v}\{N_x+\sqrt{\beta\theta^2/\alpha}\,uQ\}.
\end{aligned}                                                   \tag{IC26}
\]
Each is its leading root $P,-Q,N_x$ plus an infinitesimal-coefficient
combination of (IC22). Indeed $\theta/\sqrt\alpha=O(\pi)$,
$\beta\theta^2/\alpha=O(\pi)$, and $m,\alpha$ are infinitesimal.
Thus the complete original parent has a bounded Gram on (IC22) with
strictly positive diagonal residue
\[
             (\gamma,H_0,1,1,1,\kappa_R,\kappa_S).
                                                               \tag{IC27}
\]
This proves the reciprocal-ratio chart. Together with the bounded-ratio
chart and the inherited remote theorem it proves the stated all-tests
interior source theorem.

\paragraph{Verification and exact scope.}
The accompanying checker verifies (IC7)--(IC8), every full first
normal by original-parent differentiation, the complete hard residues,
all canceled source coquotients and contact jets, the finite Taylor
identity, and the exact remote-ratio higher-angle identity and both
strict pivots. Its denominator report specializes each canceled
source denominator at $x=0$ and records nonzero signed products of
positive-unit factors
from (IC1); no inverse opening or negative testing-ratio power is
hidden in the proof. The ordered-field bounds and the full original
remote join are the proof above. This theorem removes exactly the
source subdomain (IC1), after the original total-scale intersection;
it does not remove infinitesimal $\eta$, unbounded $\zeta$, or an
infinitesimal $r$, $1-r$, or $j_0$ opening.
